\documentclass[10pt]{article}
% Shared presentation layer for the standalone R0--R7 development documents.
% Method equations remain authoritative in the canonical idea sketch.
\usepackage[a4paper,margin=18mm,headheight=14pt]{geometry}
\usepackage{fontspec}
\usepackage{xeCJK}
\xeCJKsetup{CJKspace=true}
\setmainfont{DejaVu Serif}
\setsansfont{DejaVu Sans}
\setmonofont{DejaVu Sans Mono}
\setCJKmainfont{Noto Serif CJK KR}
\setCJKsansfont{Noto Sans CJK KR}
\setCJKmonofont{Noto Sans Mono CJK KR}

\usepackage{amsmath,amssymb}
\usepackage{booktabs,longtable,tabularx,array}
\usepackage{enumitem}
\usepackage{xcolor}
\usepackage{hyperref}
\usepackage{fancyhdr}
\usepackage{microtype}
\usepackage[most]{tcolorbox}

\definecolor{CoreBlue}{HTML}{1F4E79}
\definecolor{CoreLight}{HTML}{EAF3F8}
\definecolor{StopRed}{HTML}{9C0006}
\definecolor{StopLight}{HTML}{FCE8E6}
\definecolor{GateGreen}{HTML}{38761D}
\definecolor{GateLight}{HTML}{EAF4E2}
\definecolor{DecisionOrange}{HTML}{B45F06}
\definecolor{DecisionLight}{HTML}{FFF2CC}

\hypersetup{
  colorlinks=true,
  linkcolor=CoreBlue,
  urlcolor=CoreBlue
}

\setlength{\parindent}{0pt}
\setlength{\parskip}{0.3em}
\setlength{\emergencystretch}{3em}
\setlist[itemize]{leftmargin=1.45em,itemsep=0.1em,topsep=0.2em}
\setlist[enumerate]{leftmargin=1.75em,itemsep=0.14em,topsep=0.22em}
\renewcommand{\arraystretch}{1.15}
\newcolumntype{L}[1]{>{\raggedright\arraybackslash}p{#1}}
\newcommand{\code}[1]{{\small\nolinkurl{#1}}}
\newcommand{\must}[1]{\textcolor{CoreBlue}{\textbf{#1}}}
\newcommand{\haltitem}[1]{\textcolor{StopRed}{\textbf{#1}}}
\newcommand{\passitem}[1]{\textcolor{GateGreen}{\textbf{#1}}}
\newcommand{\decisionitem}[1]{\textcolor{DecisionOrange}{\textbf{#1}}}

\newenvironment{contractbox}[1]
  {\begin{tcolorbox}[colback=CoreLight,colframe=CoreBlue,title={#1},fonttitle=\bfseries,before upper=\raggedright]}
  {\end{tcolorbox}}
\newenvironment{decisionbox}[1]
  {\begin{tcolorbox}[colback=DecisionLight,colframe=DecisionOrange,title={#1},fonttitle=\bfseries,before upper=\raggedright]}
  {\end{tcolorbox}}
\newenvironment{stopbox}[1]
  {\begin{tcolorbox}[colback=StopLight,colframe=StopRed,title={#1},fonttitle=\bfseries,before upper=\raggedright]}
  {\end{tcolorbox}}


\hypersetup{
  pdftitle={Development Roadmap for Response-Distilled Global-Local Wind Dynamics},
  pdfauthor={Wind3DGS Project}
}

\pagestyle{fancy}
\fancyhf{}
\lhead{\small Response-Distilled Global--Local Development}
\rhead{\small Master Roadmap}
\cfoot{\thepage}

\title{\textbf{Development Roadmap}\\
Response-Distilled Global--Local Wind Dynamics}
\author{Wind3DGS Project}
\date{2026-08-22}

\begin{document}
\maketitle

\begin{contractbox}{Authority and document split}
방법 수식과 claim의 authority는
\code{3dgs_response_distilled_global_local_wind_dynamics_2026-08-22.tex}다.
이 문서는 R0--R7의 의존성, 상태와 claim routing만 소유하는 master roadmap이다.
구체적인 구현 계약, 미결정 설계, fixture와 stage 산출물은
\code{development/} 아래의 독립 문서가 각각 소유한다.
이전 단일 checklist의 세부 항목을 이 roadmap에 다시 복제하지 않는다.
\end{contractbox}

\tableofcontents
\newpage

\section{사용 목적}

개발자는 현재 작업 중인 R-stage 문서만 열어 수정한다. 이 roadmap은 다음 경우에만 갱신한다.
\begin{itemize}
  \item stage 진입/종료 상태 또는 dependency가 바뀔 때
  \item sketch stable label의 소유 stage가 바뀔 때
  \item Gate routing이나 허용 claim이 바뀔 때
  \item 새 stage가 기존 stage를 대체할 때
\end{itemize}

각 part 문서는 공통적으로 목적/소유권, 진입 입력, 동결 계약, 구현 체크리스트,
Open Design Decisions, 산출물, fixture/metric과 종료 조건을 갖는다.
설계가 끝나지 않은 항목을 구현자가 암묵적으로 선택하지 않고 해당 part의
\textbf{Open Design Decisions}에서 먼저 동결한다.

\section{문서 계층}

\begin{longtable}{L{0.22\linewidth}L{0.27\linewidth}L{0.42\linewidth}}
\toprule
계층 & 문서 & 소유권\\
\midrule
\endfirsthead
\toprule
계층 & 문서 & 소유권\\
\midrule
\endhead
Method authority &
\code{3dgs_response_distilled_global_local_wind_dynamics_2026-08-22.tex} &
연구 질문, 수식, method identity, Gate와 claim boundary\\
Master roadmap &
현재 문서 &
R0--R7 dependency/status, stable-label routing, Definition of Done\\
Part contract &
\code{development/r0_*.tex}--\code{development/r7_*.tex} &
해당 stage의 구현 계약, 미결정 설계, fixture, artifact와 exit evidence\\
R1 동반 연구 기록 &
\code{development/r1_teacher_implementation_record.tex} &
Teacher 전체 개발 과정의 수식·실패·수치·재현 경로; R1 acceptance를 대체하지 않음\\
Executable truth &
\code{../code} schema/config/test 및 \code{../experiments} manifest &
part 문서를 구현한 versioned code, 실제 threshold 값, run과 evidence\\
\bottomrule
\end{longtable}

문서와 code가 충돌하면 먼저 code를 완료 상태로 간주하지 않고 owning part와 sketch를 대조한다.
Method identity 변경이면 sketch를 먼저 수정하고, implementation-only 선택이면 owning part와
versioned config/schema를 함께 갱신한다.

\section{R0--R7 개발 문서}

\begin{longtable}{L{0.07\linewidth}L{0.33\linewidth}L{0.24\linewidth}L{0.25\linewidth}}
\toprule
Stage & Owning document & Entry & Exit\\
\midrule
\endfirsthead
\toprule
Stage & Owning document & Entry & Exit\\
\midrule
\endhead
R0 & \code{development/r0_contract_and_schema.tex} &
Canonical sketch &
Closed API/schema/registry와 no-leak contract\\
R1 & \code{development/r1_teacher_probe_oracle.tex} &
R0 artifacts &
Converged teacher, common-probe, traction/ZOH, oracle와 base transport smoke\\
R2 & \code{development/r2_single_case_global_overfit.tex} &
R1 pass &
Single-case Global overfit와 long-horizon finite rollout\\
R3 & \code{development/r3_multi_resplat_measure.tex} &
R2 pass, independent GS variants &
Gate A\(_T\)/A\(_M\)/A\(_L\) routing evidence\\
R4 & \code{development/r4_heldout_global.tex} &
R3 route 확정 &
Object-disjoint Global Gate B와 frozen Best Global\\
R5 & \code{development/r5_local_residual.tex} &
Gate B pass &
Residual atlas, oracle/always-on Local과 Gate C\\
R6 & \code{development/r6_conditional_local_runtime.tex} &
Gate C pass &
Deterministic Local state machine과 Gate D\\
R7 & \code{development/r7_renderer_and_paper_evidence.tex} &
선택한 Global/Local route &
Optional appearance, champion/failure evidence와 paper freeze\\
\bottomrule
\end{longtable}

\section{Training stage와 R-stage 대응}

\begin{longtable}{L{0.09\linewidth}L{0.22\linewidth}L{0.58\linewidth}}
\toprule
Training & Development & 동결된 의미\\
\midrule
\endfirsthead
\toprule
Training & Development & 동결된 의미\\
\midrule
\endhead
T0 & R0--R1 &
SI/schema, teacher convergence, common-probe, oracle, traction/exact-ZOH와 base transport\\
T1 & R2 &
R2 내부의 area measure 및 patch-to-token initialization; mass는 homogeneous rule로 파생\\
T2 & R2 &
Single-object/single-GS Global response overfit\\
T3 & R3--R4 &
Multi-resplat paired training과 object-disjoint held-out Global\\
T4 & R5 &
Frozen Best Global residual로 always-on Local 학습\\
T5 & R6 &
Analytic-only selector/state machine 구현, 통합과 replay 평가; selector score와 automaton은 학습하지 않음\\
T6 & R5--R6 optional &
Global field/pole과 Gate B checkpoint를 freeze한 Local/package continuity fine-tuning만 허용.
Local/package 변경은 Gate C/D를 무효화해 R5/R6을 다시 실행하고,
T6 안의 Global 변경은 금지한다. 별도 method lineage에서 Global을 바꾸면 R4/T3 Gate B부터
새 checkpoint/evidence를 만들고 이후 R5/R6을 다시 실행한다.\\
Paper & R7 &
Base transport 재사용, optional appearance와 paper evidence만 추가\\
\bottomrule
\end{longtable}

R2a/R2b는 canonical T-stage를 추가하는 이름이 아니라 R2 문서 내부의 개발 substep이다.
R2a는 T1 initialization 뒤 T2 fixed-seconds overfit을 닫고, R2b는 mandatory baseline과
schedule resolution을 닫는다. Stable teacher peak가 없으면 R2b의 period branch만 생략하고
frozen fixed-seconds fallback과 baseline 결과를 남긴다.

\section{Stable sketch reference routing}

Stable label은 sketch에만 정의하고 part 문서는 이름으로 참조한다.
\begin{longtable}{L{0.10\linewidth}L{0.36\linewidth}L{0.43\linewidth}}
\toprule
Owner & Stable label & 계약\\
\midrule
\endfirsthead
\toprule
Owner & Stable label & 계약\\
\midrule
\endhead
R0 & \code{eq:rd-setup-input}, \code{eq:rd-frame-input} &
Setup/frame I/O와 target-forbidden field\\
R0 & \code{eq:rd-typed-control-map} &
Base/effective pole, control range와 reset\\
R0 & \code{eq:rd-gaussian-token}, \code{eq:rd-latent-anchor} &
Normalization, patch-first token과 global context\\
R0 & \code{eq:rd-response-package} &
Cached package payload, schema와 identity\\
R1 & \code{eq:rd-common-probe-map}, \code{eq:rd-common-probe-adjoint} &
Training/evaluation mapping과 payload별 work adjoint\\
R1 & \code{eq:rd-teacher-consistent-mass}, \code{sec:rd-teacher-development-evidence} &
Teacher DOF/quadrature/mass와 고차 map law, 개발 증거의 채택 경계\\
R1 & \code{eq:rd-exact-zoh}, \code{eq:rd-wind-pullback} &
Frame-start traction과 single exact-ZOH evaluator\\
R1--R2 & \code{eq:rd-attachment-whitening}, \code{eq:rd-global-response} &
Exact attachment, final-coordinate pole과 passive Global\\
R2 & \code{eq:rd-response-loss}, \code{eq:rd-total-loss},
\code{eq:rd-teacher-period-spectrum} &
Common-probe response loss와 rollout schedule\\
R3 & \code{eq:rd-area-mass} &
Positive area, homogeneous mass와 density consistency\\
R5 & \code{eq:rd-local-residual}, \code{eq:rd-local-support-projection} &
Frozen-Global residual과 free/valid ordered sparse support\\
R6 & \code{eq:rd-local-score}, \code{eq:rd-local-runtime-state} &
Analytic selector와 deterministic live state\\
R6 & \code{eq:rd-mean-transport}, \code{eq:rd-pareto-gate} &
Transport-only fade와 quality--latency Gate D\\
R7 & \code{eq:rd-optional-appearance} &
Gate 후 stateless bounded appearance branch\\
\bottomrule
\end{longtable}

\section{R0 cross-cutting registries}

다음 schema의 \textbf{형식과 freeze 시점}은 R0가 소유한다. 실제 stage별 값은 owning part와
versioned config가 채운다.
\begin{itemize}
  \item \code{DomainContract}: SI unit, geometry/material/attachment/wind/OOD와 failure denominator
  \item \code{TensorAndFrameSchema}: coordinate frame, tensor shape/axis/dtype, permutation과 ragged mask
  \item \code{PhysicalValidityContract}: valid index, area/mass/load/field zeroing 또는 asset reject
  \item \code{ResponsePackageSchema}: rank, aligned payload, serialization, hash와 reset
  \item \code{MetricRegistry}: metric ID/unit/direction, mask, reduction, object/seed aggregate
  \item \code{ThresholdRegistry}: Gate Boolean expression, numeric margin, freeze split와 version
  \item \code{StageArtifactManifest}: required input hash, output/report, validator, failure code와 next route
\end{itemize}

\section{현재 상태}

\begin{longtable}{L{0.09\linewidth}L{0.20\linewidth}L{0.58\linewidth}}
\toprule
Stage & 상태 & 다음 동작\\
\midrule
\endfirsthead
\toprule
Stage & 상태 & 다음 동작\\
\midrule
\endhead
R0 & 계약 재동결 필요 & R0 Open Design Decisions와 executable schema를 먼저 닫는다.\\
R1 & 개발 진단 진행; acceptance 미완료 & Registry/샘플·시간·공간 진단 근거는 R1 문서와 동반 연구 기록에 보존. R0 잔여 계약, 최종 동일-backend 수렴과 oracle/transport를 닫는다.\\
R2 & Blocked by R1 & Network-free oracle와 base transport가 통과한 뒤 시작한다.\\
R3 & 대기 & Independent reconstruction eligibility와 density registry를 동결한다.\\
R4 & 대기 & Dataset size와 Best Global selection protocol을 동결한다.\\
R5 & 대기 & Gate B 뒤 support proposal과 residual ownership을 동결한다.\\
R6 & 대기 & Gate C 뒤 transition table과 active/live budget을 동결한다.\\
R7 & 대기 & 선택된 claim route 뒤 optional paper work만 수행한다.\\
\bottomrule
\end{longtable}

\section{Gate routing과 허용 claim}

\begin{longtable}{L{0.14\linewidth}L{0.30\linewidth}L{0.45\linewidth}}
\toprule
Gate & 실패 의미 & 허용 route\\
\midrule
\endfirsthead
\toprule
Gate & 실패 의미 & 허용 route\\
\midrule
\endhead
A\(_T\) & Teacher/map/oracle/Global prerequisite 실패 & Learned-response pipeline 중단\\
A\(_M\) & Independent-GS measure/response consistency 실패 & Fixed reconstruction family; density claim 철회\\
A\(_L\) & Latent token/relation 이득 없음 & 실패를 기록하고 capacity-matched direct-set을 R4 main route로 전달; 이후 Gate B/C 판정\\
B & Held-out Global response 실패 & Family 축소 또는 learned-response pivot 중단\\
C & Learned Local 필요성/우위 없음 & Global-only response paper\\
D & Conditional execution Pareto 실패 & Always-on Local 또는 Global-only; selector claim 철회\\
\bottomrule
\end{longtable}

R3는 claim이 축소될 수 있어도 실행 자체를 생략하지 않는다.
Total area/mass 합은 construction sanity이고, Gate A\(_M\)의 primary evidence는 spatial surface measure,
analytic traction-field force/work와 common-probe response dispersion이다.

\section{Definition of Done}

\begin{itemize}
  \item[$\square$] R0 schema/registry/no-leak contract와 serialization/permutation fixture 통과
  \item[$\square$] R1 teacher spatial/temporal convergence, common-probe/work adjoint,
  same-law traction, corrected oracle, exact-ZOH와 base mean/covariance transport 통과
  \item[$\square$] R2 single-case Global overfit와 final supported physical-time rollout 통과
  \item[$\square$] R3를 실행하고 Gate A\(_T\)를 통과하며 A\(_M\)/A\(_L\) 결과에 맞춰 claim route 기록
  \item[$\square$] R4 object-disjoint Gate B와 frozen Best Global checkpoint/report 통과
  \item[$\square$] Local을 주장하면 R5/Gate C 통과
  \item[$\square$] Conditional Local을 주장하면 R6/Gate D와 exhaustive transition replay 통과
  \item[$\square$] 선택한 route의 R7 champion/failure evidence, baseline, profiler와 reproducibility 공개 가능
  \item[$\square$] In-envelope package reject, divergence, invalid covariance와 fallback을 fixed denominator에 유지
  \item[$\square$] 이전 explicit-scaffold 결과를 새 방법의 완료 상태로 승계하지 않음
\end{itemize}

\section{업데이트 규칙}

\begin{enumerate}
  \item 진행 중인 stage의 part 문서와 code/config/test를 함께 수정한다.
  \item 새 결정은 part의 Open Design Decisions에서 frozen contract로 이동하고 version/hash를 기록한다.
  \item 다른 stage가 소비하는 schema가 바뀌면 master dependency와 downstream artifact를 함께 invalidate한다.
  \item Method equation, contribution 또는 claim boundary가 바뀔 때만 canonical sketch를 수정한다.
  \item Part TeX를 수정하면 해당 PDF를, shared preamble을 수정하면 모든 part PDF를 재빌드한다.
  \item Part/shared/master 또는 그 PDF가 바뀌면 checklist 전달 bundle도 상대 경로를 보존해 재생성한다.
  \item Master roadmap은 세부 실험 로그를 소유하지 않는다. 실제 run/evidence는 \code{../experiments}가 소유한다.
\end{enumerate}

\end{document}
